% ===========================================================================
%  Presentation
% ---------------------------------------------------------------------------
%  Same commands as the notes, same metadata, same palette. Options match
%  style/project: [mono] for black and white, [draft] for a DRAFT stripe.
%
%  Build:  latexmk -pdf slides/slides.tex        -> out/slides.pdf
% ===========================================================================
\documentclass[aspectratio=169,11pt,t]{beamer}

\makeatletter
\def\input@path{{../}{}}
\makeatother
\usepackage{style/slides}

% Basic metadata
\newcommand{\projecttitle}{Project Title}
\newcommand{\projectsubtitle}{Optional Subtitle}
\newcommand{\projectauthor}{Jonas von Stein}
\newcommand{\projectdate}{\today}
\newcommand{\projectinstitution}{Ludwig-Maximilians-Universit\"at M\"unchen}
\newcommand{\projectdepartment}{Department / Faculty}
\newcommand{\projectsupervisor}{Supervisor / Lecturer}
\newcommand{\projecttype}{Notes / Thesis / Report}
\newcommand{\projectlogo}{figures/lmu_siege1.pdf}

% Bibliography
\newcommand{\projectbibliographyfile}{ref/references}
\newcommand{\projectbibliographystyle}{alpha}

% Optional declaration block for formal theses
\newif\ifincludeauthordeclaration
\includeauthordeclarationfalse

\newcommand{\authordeclarationtitle}{Declaration of Authorship}
\newcommand{\authordeclarationtext}{%
I hereby declare that I have written this document independently and have used only the sources and aids listed in the bibliography.%
}


\begin{document}

\makeslidetitle[Seminar talk]

\begin{frame}{Outline}
    \tableofcontents
\end{frame}

% ---------------------------------------------------------------------------
%  1. Writing slides directly - every note command works unchanged
% ---------------------------------------------------------------------------
\sectionframe{Hand-written slides}

\begin{frame}{The environments you already use}
    \defn{Template idea}{
        Keep recurring setup in one place, so a document can focus on content.
    }

    \thm{Sample theorem}{
        Let \(V\) be a vector space over \(\R\). Then \(0 \cdot v = 0\) for all
        \(v \in V\).
    }
\end{frame}

\begin{frame}{Lists, math and remarks}
    \begin{itemize}
        \item Copy a paragraph straight out of \texttt{content/}.
        \item The macros are identical: \(\N \subset \Z \subset \Q \subset \R\).
        \item Permutation diagrams work too: \(\permu{1,2,3}{2,3,1}\).
    \end{itemize}

    \rmkb{
        In the box looks, boxes are unbreakable on slides. If one overflows,
        split the frame in two or allow frame breaks.
    }
\end{frame}

% ---------------------------------------------------------------------------
%  2. Importing a chapter wholesale
% ---------------------------------------------------------------------------
\sectionframe{Imported from the notes}

% \inputcontent turns \chapter and \section inside the file into slide
% headings and starts a new slide at each one.
\begin{frame}[allowframebreaks]{Foundations}
    \inputcontent{content/02_foundations}
\end{frame}

% ---------------------------------------------------------------------------
%  3. Closing
% ---------------------------------------------------------------------------
\begin{frame}[plain]
    \vfill
    \begin{center}
        {\Large Thank you.\par}
        \todo{add the outlook slide}
    \end{center}
    \vfill
\end{frame}

\end{document}
